\section*{Hallucinations}
Why do language models sometimes hallucinate-that is, make up information? At a basic level, language model training incentivizes hallucination: models are always supposed to give a guess for the next word. Viewed this way, the major challenge is how to get models to not hallucinate. Models like Claude have relatively successful (though imperfect) anti-hallucination training; they will often refuse to answer a question if they don't know the answer, rather than speculate. We wanted to understand how this works.

It turns out that, in Claude, refusal to answer is the default behavior: we find a circuit that is "on" by default and that causes the model to state that it has insufficient information to answer any given question. However, when the model is asked about something it knows well-say, the basketball player Michael Jordan-a competing feature representing "known entities" activates and inhibits this default circuit (see also this recent paper for related findings). This allows Claude to answer the question when it knows the answer. In contrast, when asked about an unknown entity ("Michael Batkin"), it declines to answer.

\begin{figure}[h]
\begin{center}
\captionsetup{labelformat=empty}
\caption{Michael Jordan → Basketball}
  \includegraphics[alt={},max width=\textwidth]{hallucination-known-entity-jordan}
\end{center}
\end{figure}

\begin{figure}[h]
\begin{center}
\captionsetup{labelformat=empty}
\caption{Michael Batkin → Can't Answer}
  \includegraphics[alt={},max width=\textwidth]{hallucination-unknown-entity-batkin}
\end{center}
\end{figure}

By intervening in the model and activating the "known answer" features (or inhibiting the "unknown name" or "can't answer" features), we're able to cause the model to hallucinate (quite consistently!) that Michael Batkin plays chess.

Sometimes, this sort of "misfire" of the "known answer" circuit happens naturally, without us intervening, resulting in a hallucination. In our paper, we show that such misfires can occur when Claude recognizes a name but doesn't know anything else about that person. In cases like this, the "known entity" feature might still activate, and then suppress the default "don't know" featurein this case incorrectly. Once the model has decided that it needs to answer the question, it proceeds to confabulate: to generate a plausible-but unfortunately untrue-response.
